% !TEX program = xelatex
% !TEX root = ../../TFM.tex
\documentclass[../../TFM.tex]{subfiles}

\begin{document}

% =====================================================
% CONCLUSIONES
% =====================================================

\chapter{Conclusiones}
\label{chap:conclusiones}
\thesisstartchap{conclusiones}

Este capítulo responde la pregunta de investigación a partir de la evidencia
presentada en el \autoref{chap:resultados}. La interpretación se limita a
asociaciones condicionales dentro de los países. Los modelos con efectos fijos,
los diagnósticos y las sensibilidades permiten evaluar la estabilidad de los
patrones estimados, pero no identifican efectos causales ni demuestran por sí
solos los mecanismos teóricos discutidos en la tesis.


\section{Síntesis de hallazgos y respuesta a la pregunta de investigación}
\label{sec:sintesis-hallazgos}

La pregunta que guía el estudio indaga cómo se asocian las rentas extractivas y
los canales productivos, institucionales y macroeconómicos con la transformación
estructural externa de las economías dependientes de recursos naturales. La
respuesta empírica se limita a asociaciones condicionales estimadas a partir de
los cambios observados dentro de cada país; no pretende explicar exhaustivamente
las diferencias históricas entre países.
En ese margen, la evidencia más estable se concentra en la asociación negativa
entre la intensidad macroeconómica de las rentas y la diversificación, y en el
papel de la estructura de la canasta exportadora para la complejidad. La
asociación agregada entre rentas y $ECI$ aparece en el modelo principal, pero es
más sensible a la muestra, las tendencias y la dinámica. Los resultados son
menos concluyentes para la moderación institucional, las condiciones
macroeconómicas y varias medidas de capacidades productivas.

En el modelo completo, $RENTS$ presenta una asociación negativa con la
complejidad económica y con la diversificación exportadora. Cuando $INST=0$, un
aumento de diez puntos porcentuales del PIB en las rentas extractivas se asocia
con una disminución aproximada de 0,098 puntos de $ECI$ y 0,035 puntos de
$DIVX$. Los p-valores convencionales son 0,023 y 0,009, respectivamente, y la
inferencia mediante \textit{Wild Cluster Bootstrap} mantiene significación al
5\% ($p_{boot}=0,040$ para $ECI$ y $p_{boot}=0,021$ para $DIVX$). La exclusión
de un país por vez conserva el signo negativo en las 49 iteraciones de ambas
ecuaciones.

La estabilidad no es idéntica en todos los ejercicios. Al incorporar tendencias
lineales específicas por país, $RENTS$ conserva una asociación negativa y
significativa con $DIVX$ ($p=0,020$), pero pierde precisión en $ECI$
($p=0,210$). En la muestra apareada de 869 observaciones, el coeficiente en
niveles deja de ser preciso para ambos resultados; al estimar primeras
diferencias, la asociación negativa reaparece para $DIVX$ ($p<0,001$), pero no
para $ECI$ ($p=0,313$). El mismo contraste se observa al incorporar un rezago de
la variable dependiente: $RENTS$ permanece negativo para $DIVX$ ($p=0,010$) y no
es significativo para $ECI$ ($p=0,163$). Por tanto, la relación con la
diversificación es más estable a través de las sensibilidades; la asociación
entre rentas y complejidad constituye un resultado del TWFE principal respaldado
por el bootstrap, pero no una regularidad uniforme ante cambios de muestra,
transformación o dinámica.

La estructura exportadora ofrece el resultado más definido para $ECI$. La
concentración ($HHI$), las exportaciones primarias no energéticas ($PEXP$) y las
exportaciones de combustibles ($FEXP$) presentan asociaciones negativas y son
conjuntamente significativas. En cambio, $PEXP$ se asocia positivamente con
$DIVX$. Esta diferencia confirma la importancia analítica de separar dos
dimensiones: una canasta puede ser menos concentrada porque incorpora una mayor
variedad de productos primarios, sin que ello implique una mayor sofisticación
o acumulación de capacidades productivas reveladas.

La calidad institucional no presenta evidencia robusta del papel moderador
previsto. El coeficiente de $INST$ es positivo y marginal al 10\% en la ecuación
principal de $ECI$, pero $RENTS\times INST$ no es significativo en ninguna de
las dos ecuaciones, tampoco con bootstrap ni en la sensibilidad dinámica. En
primeras diferencias, la interacción es positiva y significativa para $ECI$,
pero cambia de signo respecto del modelo en niveles y no se reproduce en las
demás variantes. Los efectos marginales de $RENTS$ son negativos en los cinco
percentiles observados de $INST$, aunque su precisión varía. En consecuencia,
la tesis no encuentra evidencia suficiente para afirmar que el indicador
institucional utilizado modifique sistemáticamente la asociación entre rentas y
transformación estructural externa. Esta conclusión no prueba que las
instituciones sean irrelevantes ni que una mejora institucional carezca de
efectos; delimita lo que puede inferirse con esta medición, este panel y estas
especificaciones.

Las variables $VOL$ y $RER$ no son significativas individual ni conjuntamente.
Por ello, el modelo no respalda directamente el canal macroeconómico formulado,
pero tampoco permite descartarlo: la medida de volatilidad captura una dimensión
específica de los términos de intercambio y puede no representar todos los
mecanismos fiscales, financieros o cambiarios relevantes.

En capacidades productivas, $HUMCAP$ se asocia positivamente con $DIVX$, pero no
con $ECI$, y pierde precisión al incorporar tendencias por país. $INNOV$ no
presenta una asociación adicional significativa. La evidencia, por tanto,
ofrece un respaldo limitado al papel del capital humano en la dispersión
exportadora y no autoriza una conclusión general sobre el efecto de las
capacidades en la complejidad.

Las dos variantes sectoriales completas muestran patrones complementarios. Las
rentas de minería y carbón se asocian negativamente con $ECI$
($\beta=-0,0122$, $p=0,003$), mientras que las rentas de petróleo y gas se
asocian negativamente con $DIVX$ ($\beta=-0,0047$, $p=0,001$). Al incorporar
ambos componentes en $M_{3}^{SIM}$, los patrones se conservan con coeficientes
de $-0,0136$ para minería y carbón en $ECI$ y de $-0,0048$ para petróleo y gas
en $DIVX$ ($p=0,001$ en ambos casos).

Las pruebas de estabilidad sectorial se aplican a $M_{3}^{SIM}$: los dos
resultados conservan el signo en las 49 exclusiones de países y mantienen
significación con bootstrap. Estas pruebas no se extrapolan como robustez
independiente de las variantes separadas. Además, las pruebas de igualdad de los
coeficientes directos no rechazan igualdad al 5\%. Los patrones sugieren
heterogeneidad sectorial, pero no demuestran que cada recurso opere mediante un
mecanismo exclusivo.

Finalmente, la sensibilidad de 2014 no aporta evidencia de un quiebre general.
Las pruebas conjuntas de cambio en $RENTS$, $INST$ y su interacción no rechazan
estabilidad al 5\% para $ECI$ ($p=0,343$) ni para $DIVX$ ($p=0,059$). El segundo
resultado es marginal al 10\%, pero debe interpretarse como sensibilidad
temporal en muestra completa, no como ruptura causal ni como base para dividir
el período de estimación.

En relación con las hipótesis, la evidencia puede resumirse así: la
\hyperref[hyp:h1]{primera hipótesis} recibe apoyo parcial porque los patrones son
compatibles con una relación adversa entre intensidad extractiva, estructura
exportadora y transformación externa, aunque la asociación de las rentas con
$ECI$ es sensible y el diseño no identifica directamente los mecanismos
propuestos; la \hyperref[hyp:h2]{segunda hipótesis} no recibe respaldo robusto,
pues la interacción institucional aislada en primeras diferencias para $ECI$ no
se reproduce en el modelo principal ni en las demás sensibilidades; la
\hyperref[hyp:h3]{tercera hipótesis} queda inconclusa; la
\hyperref[hyp:h4]{cuarta hipótesis} recibe apoyo limitado por el resultado de
capital humano en $DIVX$; y la \hyperref[hyp:h5]{quinta hipótesis} encuentra
patrones sectoriales sugestivos, aunque no diferencias concluyentes entre los
coeficientes directos.


\section{Implicaciones teóricas y contribuciones metodológicas}
\label{sec:implicaciones-teoricas}

La primera contribución consiste en desplazar el resultado de interés desde la
dependencia o las rentas hacia la transformación estructural externa. La
selección inicial mediante $DRES$ delimita economías con una inserción externa
dependiente de recursos, mientras que $RENTS$ mide la intensidad macroeconómica
de las rentas dentro de esa muestra. Esta separación evita tratar como
equivalentes la estructura exportadora utilizada para seleccionar países y la
variable explicativa que cambia a lo largo del tiempo.

La segunda contribución es distinguir sofisticación y diversificación. $ECI$ y
$DIVX$ no son medidas intercambiables: los signos opuestos de $PEXP$ muestran
que una expansión de la variedad exportada puede ocurrir sin un aumento de la
complejidad. Esta distinción permite evaluar con mayor precisión qué dimensión
de la transformación externa está asociada con cada canal teórico y evita usar
la palabra diversificación como sinónimo automático de \textit{upgrading}.

La tercera contribución es empírica. El estudio compara $M_1$ y $M_2$ como
especificaciones temáticas paralelas y concentra la interpretación en $M_3$,
estimado sobre una muestra común de 1.044 observaciones, 49 países y 23 años
efectivos. Esta regla evita atribuir a la incorporación de variables lo que
podría deberse a cambios de muestra. La inferencia principal se complementa con
bootstrap, exclusión de países, observaciones influyentes, efectos marginales,
tendencias nacionales y una sensibilidad temporal preespecificada. La selección
del estimador se examina mediante las comparaciones entre \textit{pooled OLS},
efectos fijos, efectos aleatorios y Mundlak; además, las primeras diferencias y
la especificación dinámica permiten evaluar qué resultados dependen de la
persistencia, la transformación temporal o la reducción de la muestra.

La cuarta contribución es de alcance interpretativo. Los diagnósticos temporales
y residuales se utilizan como información sobre el panel, no como mecanismos
automáticos para clasificar series, escoger estimadores o transformar el modelo.
Del mismo modo, la ausencia de significación se presenta como falta de evidencia
estadística suficiente y no como prueba de irrelevancia económica. Esta
jerarquía separa resultados centrales, evidencia complementaria y resultados no
concluyentes.

Los hallazgos son compatibles con partes de la literatura sobre especialización
extractiva y complejidad económica \parencite{Auty1993,Sachs1995,Hidalgo2009,Hausmann2014},
pero no constituyen una prueba general de la denominada maldición de los
recursos. Tampoco refutan la importancia de las instituciones o las capacidades.
La contribución del estudio es identificar regularidades condicionales dentro de
la muestra y mostrar cuáles permanecen bajo sensibilidades exigentes.


\section{Implicaciones para la política económica y productiva}
\label{sec:recomendaciones-politica}

Las implicaciones de política deben formularse como ámbitos que merecen
evaluación, no como recetas derivadas causalmente de los coeficientes. La
evidencia permite señalar cuatro orientaciones prudentes.

En primer lugar, el seguimiento de la transformación productiva debería utilizar
simultáneamente indicadores de sofisticación y dispersión exportadora. Una
política que aumente el número de productos exportados puede mejorar $DIVX$ sin
elevar $ECI$. Por ello, las evaluaciones de diversificación deberían considerar
también contenido tecnológico, capacidades requeridas, productividad y
encadenamientos, además del número o la participación de productos.

En segundo lugar, la asociación entre capital humano y $DIVX$ justifica evaluar
programas de formación y capacidades vinculados a la diversificación, pero el
resultado no establece qué intervención es eficaz ni demuestra un efecto sobre
la complejidad. Las políticas de educación técnica, formación avanzada o
investigación aplicada deberían someterse a evaluaciones específicas que midan
adopción tecnológica, productividad y resultados exportadores.

En tercer lugar, los patrones de minería-carbón y petróleo-gas sugieren que el
diagnóstico productivo debe considerar el tipo de recurso predominante. No
obstante, como la igualdad de los coeficientes directos no se rechaza al 5\%, los
resultados no justifican asignar automáticamente instrumentos distintos a cada
sector. Las estrategias de proveedores, procesamiento, contenido local o apoyo
a exportaciones alternativas requieren análisis de viabilidad, competencia,
costos fiscales y capacidades nacionales.

En cuarto lugar, la ausencia de evidencia robusta de moderación mediante $INST$
impide concluir que una mejora general de gobernanza neutralice por sí sola la
asociación negativa de las rentas. Este resultado tampoco demuestra que las
instituciones sean irrelevantes. De manera análoga, la ausencia de significación
de $VOL$ no autoriza a descartar la estabilización macroeconómica. Las reformas
institucionales, reglas fiscales o fondos de ahorro pueden tener fundamentos
independientes en la literatura y en objetivos de política, pero su diseño
específico no se deduce de las estimaciones de esta tesis y debe evaluarse con
evidencia adicional.


\section{Limitaciones y agenda futura de investigación}
\label{sec:limitaciones-agenda}

La principal limitación es la naturaleza observacional del diseño. Los efectos
fijos por país controlan factores invariantes y los efectos de año absorben
shocks globales comunes, pero no eliminan la causalidad inversa, la simultaneidad,
los errores de medición ni las variables omitidas que cambian en el tiempo. Por
ejemplo, la complejidad y la diversificación también pueden modificar las rentas,
las políticas públicas y las instituciones. La comparación entre
\textit{pooled OLS}, efectos fijos, efectos aleatorios y Mundlak respalda la
elección de TWFE para estimar asociaciones condicionales dentro de los países;
no constituye, sin embargo, una estrategia de identificación causal. Del mismo
modo, cambiar el estimador de la matriz de varianzas mejora la inferencia ante
determinadas formas de dependencia, pero no resuelve la endogeneidad. En ausencia
de un instrumento externo, una reforma o un shock cuasiexperimental defendible,
los coeficientes no identifican el efecto de una intervención sobre las rentas,
las instituciones o las capacidades productivas.

La segunda limitación surge de la persistencia de los resultados de
transformación estructural. Los rezagos de $ECI$ y $DIVX$ son positivos y
estadísticamente significativos en la especificación dinámica, lo que indica que
el nivel observado en un año depende en parte de su trayectoria previa. El TWFE
estático no modela explícitamente esa dependencia ni permite recuperar por sí
solo una senda de ajuste. A su vez, la inclusión de la variable dependiente
rezagada dentro de un estimador de efectos fijos puede producir sesgo de Nickell
en paneles con una dimensión temporal finita \parencite{Nickell1981}; por ello,
la especificación dinámica se interpreta como una sensibilidad acotada y no
como un reemplazo del modelo principal. Su coeficiente rezagado tampoco debe
leerse como evidencia causal. La tesis no adopta GMM porque no dispone de una
estrategia de momentos e instrumentos que sea sustantivamente defendible y que
evite una proliferación mecánica de instrumentos.

La tercera limitación corresponde al alcance de los resultados. $ECI$ y $DIVX$
describen la dimensión externa de la transformación estructural, pero no captan
directamente empleo, productividad sectorial, informalidad, valor agregado
doméstico ni encadenamientos entre firmas. Asimismo, los indicadores nacionales
pueden ocultar heterogeneidad regional, empresarial y sectorial. La interacción
entre $RENTS$ e $INST$ evalúa heterogeneidad condicional, no demuestra un
mecanismo de mediación institucional; de igual manera, la desagregación de las
rentas por sectores no identifica por sí sola los canales económicos que generan
las asociaciones estimadas.

La cuarta limitación es muestral y temporal. La selección inicial se fija con
$DRES\geq20\%$ en 1990--1995, mientras que la muestra econométrica común queda
reducida a 49 países y 23 años efectivos. El panel es desbalanceado, contiene
vacíos comunes y termina en 2021. Las especificaciones en primeras diferencias
y con variable dependiente rezagada se estiman sobre una muestra común más corta,
con 869 observaciones y 19 años efectivos entre 2003 y 2021. En consecuencia,
los resultados no se generalizan automáticamente a economías no dependientes, a
países excluidos por disponibilidad de información ni a la fase más reciente de
la transición energética.

La quinta limitación se relaciona con las mediciones. $INST$ resume cuatro
dimensiones amplias de gobernanza y puede no captar instituciones específicas
del sector extractivo, como reglas de licenciamiento, tributación, transparencia
o gestión de empresas estatales. $VOL$ mide la desviación móvil de las
variaciones del índice CTOT y no representa todas las fuentes de inestabilidad
macrofiscal. De forma similar, las rentas reales per cápita son medidas
monetarias y no dotaciones físicas.

La sexta limitación es la sensibilidad de algunos resultados a la dinámica y a
la transformación de las variables. Las pruebas de raíz unitaria ofrecen
evidencia mixta y, por tanto, los coeficientes en niveles no deben interpretarse
como relaciones de equilibrio de largo plazo. Las primeras diferencias reducen
el papel de tendencias persistentes, pero cambian el estimando hacia asociaciones
entre variaciones de corto plazo y reducen la muestra disponible. En esa
comparación, los resultados de $ECI$ son más sensibles, mientras que la relación
negativa entre rentas y $DIVX$ presenta mayor estabilidad. La asociación de
$RENTS$ con $ECI$ y varios coeficientes complementarios también pierden precisión
al incorporar tendencias lineales por país. Estas variaciones impiden presentar
todos los coeficientes del modelo base como regularidades igualmente robustas.

La agenda futura puede avanzar en cinco direcciones. Primero, extender las
series con datos comparables posteriores a 2021 para estudiar la transición
energética y los cambios recientes sin imponer anticipadamente un quiebre.
Segundo, combinar información comercial con matrices insumo-producto, empleo y
datos de firmas para observar capacidades y encadenamientos internos. Tercero,
incorporar medidas institucionales específicas de la gobernanza extractiva y
evaluar su soporte temporal dentro de los países. Cuarto, desarrollar diseños
con fuentes externas de variación, reformas escalonadas u otros eventos
institucionales defendibles que permitan responder preguntas causales. Quinto,
estudiar la dinámica con paneles continuos de mayor longitud y estimadores
dinámicos cuyas restricciones de momentos e instrumentos se definan y justifiquen
antes de la estimación. Una estrategia más compleja aplicada al mismo panel no
sustituye la necesidad de una fuente creíble de identificación.

En conclusión, la evidencia más estable del estudio es la asociación negativa
entre la intensidad de las rentas y la diversificación exportadora, junto con la
relación entre la composición de la canasta exportadora y la complejidad. La
asociación agregada entre rentas y $ECI$ es negativa en el TWFE principal y bajo
bootstrap, pero no conserva precisión de manera uniforme al introducir
tendencias nacionales, restringir la muestra, diferenciar las variables o
incorporar la dinámica. Tampoco se encuentra evidencia robusta de que el
indicador institucional utilizado modere sistemáticamente estas asociaciones;
esto no equivale a demostrar que las instituciones sean irrelevantes. Los
hallazgos aportan una base comparativa para orientar nuevas preguntas y
evaluaciones, pero describen asociaciones condicionales y no efectos causales ni
mecanismos identificados.

\thesisendchap{conclusiones}

\end{document}
